%! Graphing Quadratic Functions & Transformations — Unit 3, Lesson 3.1 — Lesson Plan
\documentclass[10pt]{article}
\usepackage{algebra2-article}
\usepackage{algebra2-boxes}
\usepackage{graphicx}   % -article does NOT load graphicx; needed for thumbnails/screenshots
\graphicspath{{images/}}
\newcommand{\TallMath}[1]{$\displaystyle #1\rule[-1.4em]{0pt}{3.2em}$}
\newcommand{\CourseName}{Algebra 2: Shepherd}
\newcommand{\SchoolYear}{2026--2027}
\newcommand{\MeetingLength}{55 minutes}
\newcommand{\UnitNumberName}{Unit 3: Quadratic Functions \quad}
\newcommand{\LessonNumberName}{Lesson 3.1: Graphing Quadratic Functions \& Transformations}

\begin{document}
\begin{center}
    {\Large\bfseries \CourseName: \SchoolYear \\
    {\small\bfseries \UnitNumberName \LessonNumberName}}
\end{center}

\begin{tcolorbox}[colback=forestbg, colframe=forest, boxrule=0.9pt, arc=2mm,
    left=2mm, right=2mm, top=1.5mm, bottom=1.5mm]
    \textbf{Primary Objective:} Students can move from an equation to a graph. Given a quadratic in any
    of its three forms --- \emph{vertex form} $f(x)=a(x-h)^2+k$, \emph{standard form} $f(x)=ax^2+bx+c$,
    or \emph{intercept (factored) form} $f(x)=a(x-p)(x-q)$ --- they can find the \emph{vertex}, the
    \emph{axis of symmetry}, the opening direction and width, the \emph{intercepts}, and the
    \emph{max/min}, then use those features to graph the parabola or match it to its graph. They read
    each form as a set of \emph{transformations} of the parent $y=x^2$: $f(x)+k$ shifts vertically,
    $f(x-h)$ shifts horizontally, and $a\,f(x)$ stretches/reflects. Where Lesson~3.0 \emph{read} a
    parabola, Lesson~3.1 \emph{builds} one from its rule.
    \\[2pt]
    \textbf{Standards (2023 VA SOL):} This lesson reactivates the Algebra~1 quadratic cluster
    \textbf{A.F.2} as Unit~3's graphing foundation: \textbf{A.F.2b} (determine the key characteristics
    --- zeros, $y$-intercept, vertex/max--min, domain and range --- from an equation or graph),
    \textbf{A.F.2c} (graph a quadratic using transformations $f(x)+k$ and $k\,f(x)$), and
    \textbf{A.F.2d} (connect the standard and factored forms to the graph). Algebra~2 extends the
    transformation menu to the horizontal shift $f(x-h)$ --- the full vertex form --- which is the same
    $f(x)+k,\,f(kx),\,f(x+k),\,k\,f(x)$ lens \textbf{A2.F.1} formalizes for every other family, applied
    here to the quadratic parent. Characteristics are named per \textbf{A2.F.2a/d}.
\end{tcolorbox}

\renewcommand{\arraystretch}{1.4}
\begin{skillbox}[Priority Ideas \& Skills]{goldbox}
\begin{tabularx}{\linewidth}{@{}X|X@{}}
\textbf{Priority Ideas \& Skills} & \textbf{Key Understandings (the ``why'')} \\[2pt]
\hline
\vspace{-6pt}
\begin{itemize}[leftmargin=1.1em, nosep, after=\vspace{2pt}]
  \item Read a parabola as \textbf{transformations} of $y=x^2$: $f(x)+k$ (up/down), $f(x-h)$ (left/right), $a\,f(x)$ (stretch/reflect).
  \item From \textbf{vertex form} $a(x-h)^2+k$: name the vertex $(h,k)$, axis $x=h$, direction \& width.
  \item From \textbf{standard form} $ax^2+bx+c$: find the axis $x=-\tfrac{b}{2a}$, the vertex, and the $y$-intercept $(0,c)$.
  \item From \textbf{intercept form} $a(x-p)(x-q)$: read the zeros $p,q$ and the axis $x=\tfrac{p+q}{2}$.
  \item Use those features to \textbf{graph} the parabola or match it to its graph.
\end{itemize}
&
\vspace{-6pt}
\begin{itemize}[leftmargin=1.1em, nosep, after=\vspace{2pt}]
  \item Each form \emph{advertises} a different feature; they are three descriptions of \emph{one} curve.
  \item In $f(x-h)$ the graph moves \emph{right} for a positive $h$ --- opposite the sign inside, the year's most common transformation slip.
  \item The axis of symmetry is the backbone: once you have it, the vertex is one substitution and a symmetric point pair follows.
  \item $-\tfrac{b}{2a}$ \emph{is} the axis because it sits halfway between symmetric points --- it is not a memorized mystery.
  \item Sign of $a$ still controls opening/max--min; $|a|$ controls width (stretch vs.\ compress).
\end{itemize}
\end{tabularx}
\end{skillbox}

\begin{skillbox}[{Vocabulary, Concepts, \& Theorems}]{sky}
\renewcommand{\arraystretch}{1.3}
\begin{tabularx}{\linewidth}{@{}l X@{}}
\textbf{Parent function} & The simplest function of a family; for quadratics it is $y=x^2$ (vertex at the origin, opens up). \\
\textbf{Vertex form} & $f(x)=a(x-h)^2+k$; the vertex is $(h,k)$ and the axis of symmetry is $x=h$. \\
\textbf{Standard form} & $f(x)=ax^2+bx+c$; the axis is $x=-\tfrac{b}{2a}$ and the $y$-intercept is $(0,c)$. \\
\textbf{Intercept (factored) form} & $f(x)=a(x-p)(x-q)$; the zeros are $x=p$ and $x=q$, and the axis is halfway between them. \\
\textbf{Vertical shift} & $f(x)+k$: slides the graph up ($k>0$) or down ($k<0$). \\
\textbf{Horizontal shift} & $f(x-h)$: slides the graph right ($h>0$) or left ($h<0$) --- opposite the sign inside. \\
\textbf{Vertical stretch / compression} & $a\,f(x)$: narrower if $|a|>1$, wider if $0<|a|<1$; reflects over the $x$-axis if $a<0$. \\
\end{tabularx}
\end{skillbox}

\begin{fixedskillbox}[Activate Prior Knowledge \& Spiral Review]{forestbg}
    The Warm-Up rehearses exactly what the graphing work needs: \textbf{(1)} the parent $y=x^2$ and its
    table of five points (the shape every transformation moves); \textbf{(2)} the Unit~2 transformation
    language ($f(x)+k$ up/down, $a\,f(x)$ stretch/reflect) applied to a single point; and \textbf{(3)}
    \emph{expanding} $(x-1)^2-4$ into $x^2-2x-3$, the algebra that links vertex form to standard form.
    Debrief item~3 aloud: the same parabola will appear all lesson in three costumes.
\end{fixedskillbox}

\begin{skillbox}[Hook]{forestbg}
    Put three equations on the board and claim they are \emph{the same parabola}:
    \[ f(x)=x^2-2x-3 \qquad f(x)=(x-1)^2-4 \qquad f(x)=(x+1)(x-3). \]
    Ask: without graphing, which one instantly tells you the \textbf{vertex}? Which one instantly tells
    you the \textbf{zeros}? Which one tells you the \textbf{$y$-intercept}? Reveal that all three graph
    to the identical curve from Lesson~3.0 (vertex $(1,-4)$, zeros $-1$ and $3$, $y$-intercept
    $(0,-3)$). Land the theme: \emph{a quadratic wears three costumes, and each one is cut to show off a
    different feature. Today we learn to read each costume and to draw the curve it describes.}
\end{skillbox}

\begin{skillbox}[Lesson]{forestbg}
\begin{multicols}{2}
\textbf{1. The parent and its four moves.} Start at $y=x^2$ (vertex origin, points
$(\pm1,1),(\pm2,4)$). Demonstrate each transformation on that point set: $f(x)+k$ shifts \emph{up} $k$;
$f(x-h)$ shifts \emph{right} $h$ (stress: the sign flips); $a\,f(x)$ makes it \emph{narrower}
($|a|>1$) or \emph{wider} ($0<|a|<1$) and \emph{flips} it if $a<0$. Every parabola is the parent after
these moves.

\textbf{2. Vertex form $a(x-h)^2+k$.} This form \emph{is} the transformations written down: vertex
$(h,k)$, axis $x=h$. Work $f(x)=(x-1)^2-4$: right~1, down~4, opens up, vertex $(1,-4)$, min value $-4$.
Then \emph{expand} to $x^2-2x-3$ to show it is the same function --- the Warm-Up payoff.

\textbf{3. Standard form $ax^2+bx+c$.} Here the vertex is hidden. Recover the axis with
$x=-\tfrac{b}{2a}$, substitute to get the vertex $y$, and read the $y$-intercept as $(0,c)$. Work
$f(x)=x^2-6x+5$: axis $x=3$, vertex $(3,-4)$, $y$-intercept $(0,5)$. Emphasize \emph{why}
$-\tfrac{b}{2a}$ works --- it is the midpoint of any symmetric pair.

\columnbreak

\textbf{4. Intercept (factored) form $a(x-p)(x-q)$.} This form advertises the \emph{zeros}: set each
factor to zero. Because the parabola is symmetric, the axis sits \emph{halfway} between them,
$x=\tfrac{p+q}{2}$, and the vertex is one substitution away. Work $f(x)=(x+1)(x-3)$: zeros $-1,3$; axis
$x=1$; vertex $(1,-4)$ --- again the same curve.

\textbf{5. From features to a graph.} With vertex, axis, direction, and one symmetric point pair, the
parabola is determined --- plot the vertex, use symmetry to place mirror points, and the shape follows.
Practice by \emph{matching} equations to pre-drawn graphs so the feature-reading is the graded move,
not freehand drawing.

\textbf{6. Bring it home.} Return to the three costumes of $x^2-2x-3$ and line up the features each one
handed you for free. Preview: Lesson~3.2 shows how to \emph{produce} the factored form (factoring) so
we can find zeros from standard form by hand.
\end{multicols}
\end{skillbox}

\begin{skillbox}[Active Monitoring]{redbox}
Circulate during the notes practice and Tier work. Watch for and correct:
\begin{itemize}[leftmargin=1.2em, nosep]
  \item the \textbf{horizontal-shift sign flip}: reading $(x-3)^2$ as a shift \emph{left} 3 instead of right~3 (and $(x+2)^2$ as right instead of left);
  \item reporting the vertex of $a(x-h)^2+k$ as $(-h,k)$ --- forgetting that $h$ is what makes the squared factor zero;
  \item computing the axis as $x=\tfrac{b}{2a}$ (dropping the minus sign) from standard form;
  \item confusing \emph{stretch} with \emph{shift}: treating the $a$ in $a x^2$ as a vertical move.
\end{itemize}
Cold-call prompts: ``Which way does $(x-4)^2$ move the parent, and how do you know?'' ``Given
$x^2+2x-3$, what is the axis --- and why is it $-\tfrac{b}{2a}$?'' ``Which form would you \emph{want}
if I asked for the zeros? for the vertex?''
\end{skillbox}

\begin{skillbox}[Group Work \& Differentiation]{redbox}
\textbf{Three Costumes, One Curve} activity (tiered; mirrors the notes).
\begin{multicols}{3}
\textbf{Tier R --- Remediate.} From vertex-form equations, name the vertex, axis, and opening direction,
then match each equation to one of three pre-drawn parabolas.

\columnbreak
\textbf{Tier A --- Approaching.} From \emph{standard} form, find the axis with $-\tfrac{b}{2a}$, compute
the vertex and $y$-intercept, complete a five-point table, and convert one equation between vertex and
standard form by expanding.

\columnbreak
\textbf{Tier E --- Extension.} Work \emph{intercept} form: read zeros and axis, then \emph{write} a
quadratic's equation from a labeled graph, and describe the full transformation sequence that carries
$y=x^2$ onto it.
\end{multicols}
\end{skillbox}

\begin{skillbox}[Individual Work \& Assessment]{redbox}
\textbf{Exit Ticket} (independent, no notes): from \textbf{vertex form} $g(x)=(x+2)^2-1$, state the
vertex, axis, opening direction, and min/max value; from \textbf{standard form} $f(x)=x^2-4x+3$, find
the axis and vertex and the $y$-intercept; and answer an \textbf{SOL-style multiple-choice} item that
turns on the horizontal-shift direction. Collect and sort into ``got it / shift-direction slip /
vertex-sign slip / $-b/2a$ slip'' to decide whether Lesson~3.2 opens with a quick re-teach.
\end{skillbox}

\begin{skillbox}[Reinforcement \& Extension]{goldbox}
\textbf{Homework:} describe transformations of $y=x^2$ and give the resulting vertex; read features from
vertex, standard, and intercept forms; match equations to pre-drawn graphs; and convert between forms.
\textbf{Extension:} a fountain-arch model given in intercept form --- interpret its zeros, vertex, and a
symmetric height in context. \textbf{Preview:} Lesson~3.2, \emph{Solving Quadratics by Factoring} ---
how to \emph{produce} the factored form so the zeros can be found from standard form by hand.
\end{skillbox}
\end{document}
